\section{Experiments}

Our experimental design tests four specific claims about coupling patterns in LLM trustworthiness dimensions. Rather than exploratory analysis, we frame each experiment as testing a hypothesis derived from our core insight: coupling exists but is sparse and model-specific.

\subsection{Experimental Questions}

We structure our evaluation around four cascading questions:

\begin{enumerate}
\item \textbf{Does coupling exist?} (h-e1) We first establish whether co-occurrence patterns between trustworthiness dimensions exceed chance levels. Without detectable coupling, subsequent analyses are moot.

\item \textbf{Is coupling difficulty-independent?} (h-m1) Detecting correlation is insufficient---we must distinguish genuine shared vulnerabilities from spurious artifacts where hard instances fail on all dimensions simultaneously. Instance difficulty is a known confound in multi-dimensional evaluation.

\item \textbf{How broad is coupling?} (h-m2) If coupling exists and is genuine, does it span many dimension pairs (indicating broad architectural bottlenecks) or few pairs (suggesting targeted vulnerability clusters)?

\item \textbf{Are coupling profiles model-specific?} (h-c1) Do different models exhibit distinct coupling patterns---fingerprints reflecting different training approaches or architectural choices?
\end{enumerate}

\subsection{Datasets}

We evaluate three model variants (GPT-4, Claude-3-Sonnet, Llama-3-70B) on 500 instances per model, with 100 instances per trustworthiness dimension (truthfulness, robustness, fairness, safety, privacy). Each instance carries binary labels (pass/fail) for all five dimensions, enabling pairwise co-occurrence analysis across 10 dimension pairs.

\textbf{Synthetic data limitation:} All Phase 4 experiments used synthetic coupling data designed to emulate benchmark structure (MultiTrust/TrustLLM remain gated). Coupling patterns were initialized with target correlation values (phi 0.35--0.40 for truthfulness-robustness and fairness-safety pairs, phi 0.15--0.25 for remaining pairs) and perturbed with 7\% noise to introduce model-specific variation. This validates measurement methodology capabilities but does not characterize real-world coupling patterns. Real LLM coupling magnitudes remain unknown pending Phase 5 production benchmark access.

\subsection{Experimental Design}

\subsubsection{h-e1: Coupling Existence Test}

\textbf{Method:} Phi coefficient analysis on $2\times2$ contingency tables (pass/fail on dimension A vs dimension B).

\textbf{Rationale:} Phi coefficient measures effect size for binary associations, interpretable on [0,1] scale where 0 indicates independence and 1 indicates perfect correlation. Unlike chi-square alone (significance test), phi quantifies coupling strength.

\textbf{Threshold:} phi $\geq$ 0.3 (medium effect size per Cohen's benchmarks) AND $p < 0.01$ (statistical significance).

\textbf{Gate criterion:} $\geq$1 model exhibits $\geq$1 significant dimension pair. This existence proof establishes coupling as a measurable phenomenon before characterizing its properties.

\subsubsection{h-m1: Difficulty-Independence Validation}

\textbf{Method:} Dual validation approach:
\begin{enumerate}
\item \textbf{Partial correlation:} Control for instance difficulty via partial phi coefficient (phi\_AB.D where D is difficulty score)
\item \textbf{Quartile stratification:} Split instances by difficulty quartiles (Q0=easy, Q3=hard) and compute phi within each quartile
\end{enumerate}

\textbf{Rationale:} Instance difficulty is a known PMC (point of maximum correlation) confound---if hard instances fail on all dimensions, raw coupling estimates inflate. Partial correlation isolates genuine coupling from this artifact. Quartile stratification provides complementary non-parametric validation: if coupling persists across difficulty levels, it is not driven solely by hard instances.

\textbf{Threshold:} partial phi $\geq$ 0.25 for $\geq$2 dimension pairs AND coupling observable in $\geq$3 of 4 difficulty quartiles.

\textbf{Gate criterion:} MUST\_WORK. Difficulty-independence distinguishes our claim (shared vulnerability mechanisms) from alternative explanation (measurement artifacts).

\subsubsection{h-m2: Coupling Breadth Analysis}

\textbf{Method:} Count significant dimension pairs per model with Bonferroni correction for multiple testing (alpha = 0.01/30 tests $\approx$ 0.00033).

\textbf{Rationale:} Testing 10 dimension pairs $\times$ 3 models = 30 comparisons inflates family-wise error rate. Bonferroni correction controls false positives at cost of conservative thresholds. We preregistered $\geq$3 pairs per model as ``broad coupling'' threshold based on combinatorial argument: 10 pairs = sparse (1--2 pairs), moderate (3--5 pairs), broad (6+ pairs).

\textbf{Threshold:} $\geq$3 significant pairs (phi $\geq$ 0.3, p\_adj $< 0.01$) for $\geq$2 of 3 models.

\textbf{Gate criterion:} SHOULD\_WORK (exploratory). This tests our sparse coupling hypothesis (FAIL = narrow coupling, PASS = broad coupling).

\subsubsection{h-c1: Model-Specific Fingerprints}

\textbf{Method:} Mantel test comparing $10\times10$ coupling matrices (phi values for all dimension pairs) across model pairs. Permutation test (10,000 iterations) generates null distribution of matrix correlations.

\textbf{Rationale:} If coupling profiles are model-specific, coupling matrices should differ structurally (low Pearson $r$ between matrices). Mantel test accounts for matrix structure and provides significance test via permutation.

\textbf{Threshold:} Mantel $r < 0.7$ (low similarity) AND $p < 0.0167$ (Bonferroni-corrected for 3 model pairs) for $\geq$1 model pair.

\textbf{Gate criterion:} SHOULD\_WORK. Model-specificity is the deployment-relevant finding (coupling as selection criterion).

\subsection{Metrics}

\begin{itemize}
\item \textbf{Phi coefficient:} Effect size for $2\times2$ contingency tables, computed as sqrt(chi-square / n). Ranges [0,1] with interpretable thresholds (0.1=small, 0.3=medium, 0.5=large).
\item \textbf{Chi-square p-value:} Statistical significance test for independence hypothesis. Alpha = 0.01 (conservative) to control false positives.
\item \textbf{Partial phi:} Phi coefficient after removing linear relationship with difficulty score. Implemented via pingouin.partial\_corr with Pearson method.
\item \textbf{Mantel r:} Pearson correlation between vectorized upper-triangle entries of two coupling matrices. Permutation-based p-value via scikit-bio.stats.distance.mantel.
\end{itemize}

\subsection{Statistical Considerations}

\textbf{Multiple testing correction:} h-m2 applies Bonferroni correction (alpha\_adj = 0.01/30). h-e1/h-m1 test prespecified pairs (truthfulness-robustness, fairness-safety) without correction, justified by prior work documenting these couplings.

\textbf{Power analysis:} Mantel test requires $n \geq 500$ instances/dimension for detecting $r = 0.5$ at 80\% power (Legendre \& Legendre). Our sample (n = 100/dimension) is underpowered for h-c1, acknowledged as limitation.

\textbf{Difficulty independence assumption:} Synthetic difficulty scores generated as Normal(0.5, 0.15) with constraint $|\text{corr}(\text{difficulty}, \text{dimension})| < 0.2$. This ensures difficulty is uncorrelated with dimensions, validating independence assumption for partial correlation.

\subsection{Why These Choices Test the Claims}

Each experiment maps directly to a claim in our sparse coupling hypothesis:

\begin{itemize}
\item \textbf{h-e1 $\rightarrow$ ``Coupling exists'':} Phi coefficient quantifies co-occurrence strength; significance test rejects independence hypothesis.
\item \textbf{h-m1 $\rightarrow$ ``Coupling reflects shared vulnerabilities'':} Partial correlation isolates genuine coupling from difficulty confound; quartile persistence shows coupling isn't driven by hard-instance artifacts.
\item \textbf{h-m2 $\rightarrow$ ``Coupling is sparse'':} Counting significant pairs across models reveals whether coupling is narrow (1--2 pairs) or broad (3+ pairs).
\item \textbf{h-c1 $\rightarrow$ ``Coupling profiles are model-specific'':} Mantel test detects structural differences in coupling matrices, indicating distinct vulnerability architectures.
\end{itemize}

This cascading design builds evidence incrementally: existence $\rightarrow$ genuineness $\rightarrow$ breadth $\rightarrow$ specificity. Each experiment's success strengthens subsequent tests' interpretability.
